\section{Related Work}

\paragraph{Driving reconstruction and sensor returns.}
3D Gaussian Splatting provides an efficient appearance representation~\cite{kerbl2023gaussians}.
UniSim, LiDAR4D, and DyNFL extend neural reconstruction to driving scenes and time-varying sensor observations
\cite{yang2023unisim,zheng2024lidar4d,wu2024dynfl}. LiDAR additionally measures free-space intervals before
its endpoints~\cite{hu2020visibility}. Neural LiDAR Fields learns a ray-weight distribution for returns
\cite{huang2023nlf}, and CaDDN represents metric depth categorically~\cite{reading2021caddn}.
Our focus is the relationship between an Actor's completed surface and the return distribution induced by
that surface. We evaluate these two objects separately and connect them through primitive-level evidence.

\paragraph{Completion and evidential geometry.}
Object-centric occupancy completion combines canonicalized LiDAR tracklets with implicit decoding
\cite{zheng2024objectcentric}. PoinTr learns point-set completion, while SnowflakeNet recursively grows
child points~\cite{yu2021pointr,xiang2021snowflakenet}. Gau-Occ uses completed LiDAR geometry to initialize
a Gaussian occupancy representation~\cite{lv2026gauocc}. These approaches motivate our target-set construction
and child-surface decoder. EAS further supervises canonical geometry with frame coverage and measured ray depth.
ALSO and evidential occupancy methods use sensor observations to distinguish free, occupied, and unobserved
space~\cite{boulch2023also,kalble2024accurate,kalble2025evocc}. GaussianFormer-2, GaussRender, and Vol3DGS
relate Gaussian support to occupancy or ray transmittance
\cite{huang2025gaussianformer2,chambon2025gaussrender,talegaonkar2025vol3dgs}.
We study how continuous evidence should compose on a fixed surface, comparing additive transmittance with
categorical return mass.

\paragraph{Motion and appearance factorization.}
Motion compensation is essential when constructing targets from accumulated sweeps~\cite{xia2023scpnet}.
DualAD separates static and dynamic streams~\cite{doll2024dualad}; DynamicVGGT predicts current and future
point maps and supervises Gaussian motion with scene flow~\cite{he2026dynamicvggt}.
Neural scene graphs provide object transforms for dynamic reconstruction~\cite{ost2021neural}.
EAS uses annotated rigid motion to canonicalize evidence and relocate the learned surface.
Its visual state follows the attribute-on-geometry principle explored by Feature 3DGS
\cite{zhou2024feature3dgs}, with a stop-gradient physical carrier that gives appearance updates a separate owner.
